% !TeX root = ../DECO_buku_iot_cerdas.tex
\chapter{INTEGRASI KECERDASAN BUATAN}

\section{Konsep AI pada Sistem}

Kecerdasan buatan pada sistem ini digunakan untuk memprediksi suhu 1 menit
ke depan berdasarkan data historis suhu yang tersimpan di ThingSpeak. AI tidak
dijalankan secara langsung pada ESP32-S3, melainkan dijalankan pada komputer
menggunakan Python. Pendekatan ini dipilih agar proses pelatihan dan pengujian
model lebih sederhana, mudah dimodifikasi, dan sesuai untuk tugas berbasis
simulasi.

Model AI yang digunakan adalah \textit{Random Forest Regressor}. Random
Forest merupakan metode \textit{ensemble learning} yang membangun sejumlah
pohon keputusan dan menggabungkan hasil prediksi dari pohon-pohon tersebut
\cite{breiman2001}. Dalam implementasi scikit-learn, \textit{RandomForestRegressor}
bekerja dengan membangun beberapa \textit{decision tree regressor} pada
\textit{subset} data dan menggunakan rata-rata hasil prediksi untuk
meningkatkan akurasi serta mengendalikan \textit{overfitting}
\cite{sklearn_rf}.

\section{Dataset dan Data Historis}

Dataset yang digunakan berasal dari data suhu DS18B20 yang dikirim oleh
ESP32-S3 ke ThingSpeak. Data dikirim setiap 15 detik (batas aman minimum
ThingSpeak). Karena target prediksi adalah 1 menit ke depan, maka model
membutuhkan target suhu 4 data setelah data saat ini. Perhitungan jumlah
langkah prediksi adalah:

\begin{equation}
\text{Jumlah Langkah} = \frac{60 \text{ detik}}{15 \text{ detik/data}} = 4 \text{ data}
\label{eq:pred-steps}
\end{equation}

Channel ThingSpeak yang digunakan bersifat \textit{public}, sehingga
pembacaan data tidak memerlukan \textit{Read API Key}. Dengan demikian,
model belajar dari pola data saat ini dan data sebelumnya untuk memprediksi
suhu 4 langkah ke depan. Pendekatan berbasis data historis
ini sesuai dengan konsep regresi deret waktu sederhana, karena input model
dibentuk dari nilai suhu aktual, nilai suhu sebelumnya, perubahan suhu, dan
rata-rata bergerak.

\section{Metode Random Forest Regressor}

\textit{Random Forest Regressor} digunakan karena mampu menangani hubungan
data yang tidak sepenuhnya linear dan relatif stabil terhadap variasi data.
Model ini tidak hanya menggunakan satu nilai suhu aktual, tetapi juga
beberapa fitur turunan dari data historis, seperti suhu sebelumnya, perubahan
suhu, dan rata-rata bergerak. Secara umum, Random Forest memiliki keunggulan
karena memanfaatkan banyak pohon keputusan sehingga prediksi tidak bergantung
pada satu model tunggal \cite{breiman2001}, \cite{sklearn_rf}.

Parameter model yang digunakan dalam implementasi adalah:

\begin{enumerate}
    \item \texttt{n\_estimators = 100}: jumlah pohon keputusan dalam
    \textit{forest}.
    \item \texttt{max\_depth = 6}: kedalaman maksimum setiap pohon untuk
    mencegah \textit{overfitting}.
    \item \texttt{random\_state = 42}: nilai \textit{seed} untuk
    reproduktibilitas hasil.
\end{enumerate}

Model dilatih ulang secara periodik menggunakan data terbaru dari ThingSpeak.
Setelah model dilatih, input terbaru digunakan untuk menghasilkan prediksi
suhu 1 menit ke depan. Hasil prediksi tersebut kemudian diklasifikasikan
menjadi status dingin, normal, atau panas.

\section{Fitur Masukan Model}

Fitur yang digunakan dalam model \textit{Random Forest Regressor} ditunjukkan
pada Tabel~\ref{tab:fitur}.

\begin{table}[H]
\centering
\caption{Fitur Masukan Random Forest Regressor}
\label{tab:fitur}
\begin{tabularx}{\textwidth}{c L{3.2cm} L{8.5cm}}
\toprule
\textbf{No} & \textbf{Fitur} & \textbf{Deskripsi} \\
\midrule
1 & \texttt{temperature}   & Suhu aktual terbaru \\
2 & \texttt{temp\_lag\_1}  & Suhu satu data sebelumnya (15 detik lalu) \\
3 & \texttt{temp\_lag\_2}  & Suhu dua data sebelumnya (30 detik lalu) \\
4 & \texttt{temp\_lag\_3}  & Suhu tiga data sebelumnya (45 detik lalu) \\
5 & \texttt{temp\_change}  & Selisih suhu aktual dengan suhu sebelumnya \\
6 & \texttt{moving\_avg\_3} & Rata-rata suhu dari tiga data terakhir (45 detik) \\
7 & \texttt{moving\_avg\_5} & Rata-rata suhu dari lima data terakhir (75 detik) \\
\bottomrule
\end{tabularx}
\end{table}

\section{Klasifikasi Status Suhu}

Setelah suhu 1 menit ke depan diprediksi, nilai tersebut diklasifikasikan ke
dalam tiga kategori. Klasifikasi dilakukan menggunakan batas ambang yang telah
ditentukan. Status suhu dingin diberikan apabila suhu lebih kecil dari
25\textdegree C, status normal diberikan apabila suhu berada pada rentang
25\textdegree C hingga 30\textdegree C, dan status panas diberikan apabila
suhu lebih besar dari 30\textdegree C.

\begin{table}[H]
\centering
\caption{Kategori Status Suhu}
\label{tab:status-suhu}
\begin{tabularx}{\textwidth}{L{5cm} c c}
\toprule
\textbf{Rentang Suhu} & \textbf{Status} & \textbf{Kode AI} \\
\midrule
Suhu $< 25$\textdegree C                              & Dingin & 0 \\
$25$\textdegree C $\leq$ Suhu $\leq 30$\textdegree C & Normal & 1 \\
Suhu $> 30$\textdegree C                              & Panas  & 2 \\
\bottomrule
\end{tabularx}
\end{table}

\section{Evaluasi Model}

Evaluasi model dilakukan menggunakan metrik regresi, yaitu \textit{Mean
Absolute Error} (MAE), \textit{Root Mean Squared Error} (RMSE), dan R$^2$
\textit{Score}. MAE digunakan untuk menghitung rata-rata selisih absolut
antara nilai aktual dan nilai prediksi. RMSE digunakan untuk memberikan
penalti lebih besar terhadap \textit{error} yang tinggi. R$^2$
\textit{Score} digunakan untuk melihat kemampuan model dalam menjelaskan
variasi data target. Metrik-metrik tersebut umum digunakan pada evaluasi
model regresi dan tersedia pada modul \texttt{sklearn.metrics}
\cite{sklearn_mae}, \cite{sklearn_r2}.

Persamaan untuk masing-masing metrik adalah:

\begin{equation}
\text{MAE} = \frac{1}{n} \sum_{i=1}^{n} |y_i - \hat{y}_i|
\label{eq:mae}
\end{equation}

\begin{equation}
\text{RMSE} = \sqrt{\frac{1}{n} \sum_{i=1}^{n} (y_i - \hat{y}_i)^2}
\label{eq:rmse}
\end{equation}

\begin{equation}
R^2 = 1 - \frac{\sum_{i=1}^{n}(y_i - \hat{y}_i)^2}
               {\sum_{i=1}^{n}(y_i - \bar{y})^2}
\label{eq:r2}
\end{equation}

di mana $y_i$ adalah nilai aktual, $\hat{y}_i$ adalah nilai prediksi, dan
$\bar{y}$ adalah rata-rata nilai aktual.

\begin{table}[H]
\centering
\caption{Metrik Evaluasi Model}
\label{tab:metrik}
\begin{tabularx}{\textwidth}{l X X}
\toprule
\textbf{Metrik} & \textbf{Fungsi} & \textbf{Interpretasi} \\
\midrule
MAE         & Mengukur rata-rata \textit{error} absolut antara nilai prediksi dan nilai aktual         & Semakin kecil nilainya semakin baik performa model \\
RMSE        & Mengukur akar rata-rata kuadrat selisih antara nilai prediksi dan nilai aktual           & Semakin kecil nilainya semakin baik; lebih sensitif terhadap \textit{error} besar \\
R$^2$ Score & Mengukur proporsi variasi data target yang dapat dijelaskan oleh model                  & Semakin mendekati 1,0 semakin baik; nilai negatif berarti model lebih buruk dari rata-rata \\
\bottomrule
\end{tabularx}
\end{table}

Data pelatihan dan pengujian dibagi menggunakan rasio 80:20 dari total data
yang tersedia. Rasio ini digunakan agar model mendapatkan cukup data untuk
belajar pola suhu, sekaligus menyisakan data yang cukup untuk mengevaluasi
kemampuan generalisasi model.

\section{Alur Kerja AI}

Alur kerja AI pada sistem adalah sebagai berikut:

\begin{enumerate}
    \item Python membaca 200 data terbaru dari ThingSpeak Field~1 melalui
    \textit{public channel} (tanpa \textit{Read API Key}).
    \item Data suhu dikonversi menjadi \textit{dataframe} menggunakan pandas.
    \item Tujuh fitur \textit{time series} dibentuk dari data historis.
    \item Target prediksi dibuat dengan menggeser data suhu sebanyak 4
    langkah ke depan (sesuai Persamaan~\ref{eq:pred-steps}).
    \item \textit{Random Forest Regressor} dilatih menggunakan 80\% data;
    jika data terlalu sedikit, seluruh data digunakan untuk \textit{training}.
    \item Evaluasi dilakukan pada 20\% data sisa menggunakan MAE, RMSE, dan
    R$^2$ \textit{Score}.
    \item Model menghasilkan prediksi suhu 1 menit ke depan menggunakan
    fitur dari data terbaru.
    \item Hasil prediksi diklasifikasikan menjadi dingin, normal, atau panas.
    \item Hasil prediksi dan status dikirim ke ThingSpeak Field~2 dan Field~3,
    dengan jeda antar pengiriman minimal 45 detik untuk menghindari
    \textit{rate limit} ThingSpeak.
    \item Proses pembacaan diulang setiap 10 detik.
\end{enumerate}

\begin{figure}[H]
\centering
\begin{tikzpicture}[
    node distance=0.45cm and 0.5cm,
    block/.style={rectangle, draw, rounded corners=3pt,
                  minimum height=0.75cm, minimum width=5.0cm,
                  text centered, font=\small, fill=blue!10},
    arrow/.style={-Latex, thick}
]
\node[block] (s1) {Baca 200 data ThingSpeak Field 1};
\node[block, below=of s1] (s2) {Konversi ke DataFrame (pandas)};
\node[block, below=of s2] (s3) {Bentuk 7 fitur \textit{time series}};
\node[block, below=of s3] (s4) {Buat target: geser 4 langkah ke depan};
\node[block, below=of s4] (s5) {Split 80/20: training \& testing};
\node[block, below=of s5] (s6) {Training Random Forest Regressor};
\node[block, below=of s6] (s7) {Evaluasi: MAE, RMSE, R$^2$ Score};
\node[block, below=of s7] (s8) {Prediksi suhu 1 menit ke depan};
\node[block, below=of s8] (s9) {Klasifikasi: Dingin / Normal / Panas};
\node[block, below=of s9] (s10){Kirim ke ThingSpeak Field 2 \& 3};
\node[block, below=of s10,fill=green!10] (s11){Tunggu 10 detik, ulangi};

\foreach \a/\b in {s1/s2, s2/s3, s3/s4, s4/s5, s5/s6, s6/s7,
                   s7/s8, s8/s9, s9/s10, s10/s11}
    \draw[arrow] (\a) -- (\b);
\end{tikzpicture}
\caption{Alur Integrasi AI Random Forest Regressor dengan ThingSpeak}
\label{fig:alur-ai}
\end{figure}
